\documentclass[conference]{IEEEtran}

\usepackage[T1]{fontenc}
\usepackage[utf8]{inputenc}
\usepackage{cite}
\usepackage{amsmath,amssymb}
\usepackage{array}
\usepackage{booktabs}
\usepackage{graphicx}
\usepackage[hidelinks]{hyperref}
\setlength{\tabcolsep}{1.5pt}

\newcommand{\verified}{\textsc{verified preserved}}
\newcommand{\partialresult}{\textsc{observed partial}}
\newcommand{\alteredresult}{\textsc{altered}}
\newcommand{\lostresult}{\textsc{confirmed lost}}
\newcommand{\unresolvedresult}{\textsc{unresolved equivalence}}
\newcommand{\measurementresult}{\textsc{measurement error}}

\title{A11yPreserve: Source-Grounded Testing of Accessibility Preservation in DOCX-to-PDF Conversion}
\author{Anonymous for review}

\begin{document}
\maketitle

\begin{abstract}
Author-supplied accessibility information in a DOCX document can be changed, removed, or represented differently during conversion even when the destination PDF is tagged. We present A11yPreserve, a source-grounded testing protocol that freezes feature contracts, verifies source OOXML, converts identical documents through named real pipelines, extracts multiple forms of PDF evidence, and compares cross-format semantics with explicit abstention. The prospectively frozen V4 primary benchmark contains 13 accessibility feature families, five controlled variants per family, 65 atomic DOCX fixtures, and 130 fixture--pipeline cases across LibreOffice and Google Docs. LibreOffice produced 38 verified-preserved, 11 observed-partial, 9 altered, 1 confirmed-lost, 1 unresolved, and 5 measurement-error atomic outcomes. Google Docs produced 11, 32, 11, 6, 5, and 0 respectively. A post-freeze supplementary Microsoft Word evaluation over the same 65 fixtures produced 16, 26, 16, 2, 0, and 5 respectively; these 65 observations are not part of the primary denominator. V5 adds corrected fixture-level holdout accounting, a documented LibreOffice profile comparison, source and destination mutation testing, independent technical validation, and calibration ground-truth auditing. The findings are descriptive of the recorded conditions: the contribution is an auditable testing methodology for selected machine-verifiable semantics, not a complete accessibility assessment or universal converter ranking.
\end{abstract}

\begin{IEEEkeywords}
document conversion, accessibility preservation, DOCX, PDF, source contracts, differential testing, testing oracles
\end{IEEEkeywords}

% The V4 sections are retained as the frozen full-paper evidence base. V5's
% completed hardening analyses are reported in the section below.
\section{Introduction}
\label{sec:introduction}

Accessibility information is authored in document representations, not only in visible appearance. A heading carries hierarchy, an image can carry an author-written alternative, a run can identify a language change, and a table can encode header relationships. These properties are often transformed before a document reaches its readers. A tagged PDF can therefore be structurally accessible while still changing or omitting a source property that an author deliberately supplied.

This distinction is recognized by accessibility guidance. ATAG 2.0 addresses preservation during restructuring and recoding transformations when equivalent output mechanisms exist \cite{atag20}. The Revised Section 508 standards and EN 301 549 likewise motivate preservation of information provided for accessibility during format conversion \cite{section508,etsi301549}. The testing problem is not simply whether the destination has tags. It is whether a known source semantic survives a representation-changing transformation, and whether the available evidence is strong enough to claim that it did.

Prior work studies accessible document conversion, PDF-to-HTML reconstruction, PDF remediation, and destination accessibility evaluation. Roig and Ribera examined office-document conversion to EPUB \cite{roig2015}; SciA11y reconstructs accessible HTML from scientific papers \cite{wang2021scia11y}; iTagPDF automates PDF accessibility metadata generation and improvement \cite{mowar2026}; and recent benchmarks study PDF accessibility evaluation and conversion services \cite{kumar2025,daisy2026,docaccessible2026}. A11yPreserve is narrower. It begins with controlled DOCX documents whose target properties are verified in the source, evaluates named DOCX-to-PDF pipelines, and separates demonstrated preservation from partial evidence, changed author values, confirmed absence, and unresolved equivalence. This is a bounded contribution to testing, not a claim that accessibility-conversion research is new.

We ask four questions:

\begin{itemize}
\item \textbf{RQ1:} To what extent do the two prospectively frozen primary pipelines preserve selected source accessibility semantics across the controlled V4 corpus?
\item \textbf{RQ2:} Which feature families exhibit verified preservation, partial preservation, alteration, confirmed loss, unresolved equivalence, or measurement difficulty?
\item \textbf{RQ3:} How do preservation outcomes differ across the two primary pipelines, and what additional feature-level observations arise in the supplementary Word arm?
\item \textbf{RQ4:} How does the source-grounded procedure behave under calibration, holdout evaluation, repeated conversion, and realistic multi-feature documents?
\end{itemize}

The contributions are fivefold: (1) a source-contract preservation methodology for controlled DOCX accessibility properties; (2) an uncertainty-aware cross-format oracle that abstains when equivalence is not established; (3) a prospectively frozen V4 benchmark with 13 families and five variants per family; (4) a primary empirical evaluation across two real DOCX-to-PDF pipelines plus a separately frozen full-corpus supplementary Word evaluation; and (5) calibration, source-oracle mutation, holdout, repeatability, and realistic-document evidence that bounds the reliability claims.

The primary unit is a fixture--pipeline case, not a document-level accessibility score. The prospectively frozen primary experiment contains 130 cases; the later Word arm adds 65 supplementary observations, for 195 archived atomic fixture--pipeline outcomes in total. The results are descriptive of the recorded conditions and do not estimate prevalence, establish user harm, or rank converters globally.

\begin{figure*}[t]
\centering
\fbox{\begin{minipage}{0.93\textwidth}
\centering
\textbf{Verified DOCX + source contract} $\longrightarrow$ \textbf{source/destination extraction} $\longrightarrow$ \textbf{feature comparison} $\longrightarrow$ \textbf{classification or abstention} $\longrightarrow$ \textbf{caveat and provenance record}
\end{minipage}}
\caption{Evidence architecture. A11yPreserve moves from a verified source contract through extraction and feature-specific comparison to a classification or abstention, while retaining the caveat and provenance record. Destination-only accessibility status is not treated as preservation evidence.}
\label{fig:workflow}
\end{figure*}

\section{Background and Related Work}
\label{sec:background}

\subsection{Preservation as a transformation requirement}
ATAG 2.0's transformation guidance treats restructuring and recoding as accessibility-relevant operations: when the destination technology provides an equivalent mechanism, accessibility information should be preserved or the author should be warned or prompted to check it \cite{atag20}. Revised Section 508 identifies preservation of information provided for accessibility in format conversion, and EN 301 549 contains a related transformation requirement \cite{section508,etsi301549}. These sources motivate the question but do not predict the feature-specific outcomes of the pipelines evaluated here.

\subsection{Office-document conversion}
Roig and Ribera studied accessibility consequences of office-document conversion to EPUB \cite{roig2015}. Their work is a direct precedent for treating conversion as an accessibility question rather than a purely syntactic operation. A11yPreserve differs in destination format, source instrumentation, and decision boundary: it uses frozen DOCX contracts, raw OOXML checks, named DOCX-to-PDF pipelines, PDF structure extraction, provenance records, and an explicit unresolved category. The distinction is one of scope and method, not a claim that prior work ignored preservation.

\subsection{Oracle construction and transformation testing}
The oracle problem is a long-standing software-testing concern: executing a test does not guarantee that the output can be judged correct \cite{barr2015oracle}. A11yPreserve uses a specification-derived feature oracle. The source contract states the intended property, while the destination comparator must determine whether a PDF representation is accepted as equivalent. Differential testing supplies the source/destination pairing, but the comparison is not a majority vote across converters. Metamorphic testing is a related strategy for constructing expected relations \cite{segura2016metamorphic}; it is not the primary oracle here because the source contracts provide explicit expected properties. Format-conversion scholarship likewise emphasizes that successful conversion requires a declared notion of what must be preserved \cite{sperberg2011conversion}.

\subsection{Accessible destination generation and evaluation}
SciA11y converts scientific PDFs to accessible HTML and focuses on navigation-oriented destination generation \cite{wang2021scia11y}. iTagPDF uses source cues and PDF analysis to improve tagging, reading order, alternative text, and formula structure \cite{mowar2026}. Kumar et al. benchmark automated and LLM-based PDF accessibility evaluation across criteria including alternative text, reading order, tables, and links \cite{kumar2025}. DAISY's benchmark compares PDF-to-accessible-document services using shared criteria and consensus scoring \cite{daisy2026}. DocAccessible reports source-grounded PDF-to-HTML benchmarking and practical Google Docs export guidance \cite{docaccessible2026,docaccessible-google}. These efforts provide important context, but they answer different questions or use different directions of conversion.

A11yPreserve's boundary is therefore modest: it does not invent accessibility-conversion benchmarking, replace destination-only validators, or claim a complete accessibility checker. It supplies a controlled source-to-destination preservation protocol whose evidence boundary is explicit and whose uncertainty is retained.

\section{A11yPreserve Method}
\label{sec:method}

A11yPreserve treats accessibility preservation as a source-to-destination comparison problem. An exact frozen DOCX is converted by a named system under test. Source evidence establishes a contract for one accessibility property; destination extraction supplies observations; and a comparison rule assigns a classification only when the evidence supports it. This framing is narrower than asking whether a PDF is accessible in every respect.

\subsection{Source contracts}
A feature contract is a machine-readable statement of the source property, its expected representation, the evidence needed to verify the source, the destination structures that can express the property, and the comparison boundary. Contracts are feature-specific because a heading hierarchy, an image alternative, a language span, a footnote association, and a table-header relation have different source and destination representations. The contract is not a weighted accessibility score and does not assert that all feature families have equal importance.

\subsection{Evidence architecture and measurement boundary}
The evidence architecture is deliberately linear: a verified source contract is paired with source and destination extraction; a feature-specific comparator then assigns a classification or abstains; the record retains the missing-evidence caveat and provenance. The protocol measures semantic preservation of selected contracted properties. It does not measure byte-for-byte serialization identity, PDF/A or PDF/UA conformance, or human usability. Those are separate questions: serialization identity concerns output encoding, conformance requires a standards-based validator, and usability requires human task or user evaluation. A tagged PDF is therefore evidence about destination structure only, not proof of any of the latter properties.

\subsection{Source verification and oracle paths}
The source-verification chain uses the frozen source manifest, raw OOXML inspection, and a source accessibility intermediate representation where applicable. Raw checks inspect the DOCX ZIP/XML parts relevant to the contract: document paragraphs and styles, drawing descriptions and relationships, numbering, language properties, tables, metadata, footnotes, equations, captions, and links. A separately implemented raw-OOXML oracle performs fresh ZIP/XML assertions without importing fixture-generation, manifest-generation, ASIR, destination-extraction, or comparator code. It detected all 39 source-oracle mutations used in the final audit.

The oracle path is a stronger cross-check against extractor-only mistakes, not external validation. It was implemented and reviewed by the same research team, and no disabled-user panel or independent external expert adjudication is claimed. The source files and manifest hashes were frozen before the final V4 conversion results.

\subsection{Conversion systems under test}
Both primary systems receive the identical frozen DOCX source and produce a tagged PDF. The LibreOffice condition used LibreOffice 26.2.6.3, Writer's \texttt{writer\_pdf\_Export}, \texttt{UseTaggedPDF=true}, \texttt{PDFUACompliance=true}, \texttt{SelectPdfVersion=1}, and an isolated profile. The Google Docs condition used an authenticated Chrome 154.0.8037.97 session, exact DOCX upload, import completion, and the visible File $\rightarrow$ Download $\rightarrow$ PDF Document (.pdf) action. Each primary conversion manifest contains 75 records: 65 atomic fixture records plus 10 project-created realistic integrated-document records. All 75 records in each manifest use the recorded UI-download route. Google output records observed the producer \texttt{Skia/PDF m156 Google Docs Renderer}; its internal service version cannot be pinned.

No content editing, print-to-PDF, screenshot conversion, post-export repair, or V3 PDF relabeling was used. Microsoft Word was evaluated separately after the primary freeze. Its native desktop GUI export used the same frozen DOCX inputs and produced valid tagged PDFs; Word is reported as a supplementary third pipeline rather than as a third prospective primary condition and is not directly comparable as a prospective treatment. External COM automation still failed with 0x80070520, but that interface limitation did not change the converter under test or the GUI export route.

\subsection{Destination evidence}
Each output record retains source and output SHA-256 values, converter metadata, timestamps, output size, PDF signature, parseability, page count, tagging observation, producer, route, and error status. Destination extraction traverses structure-tree roles and relevant objects, including headings, figures and alternative text, list containers, table and header roles, document and span language, link annotations, metadata, notes, formulas, captions, and raw object markers. pypdf traversal is supplemented by PyMuPDF text/link evidence, raw serialized-object inspection, and Poppler text extraction for difficult cases. These paths reduce but do not eliminate parser uncertainty.

\subsection{Evidence-aware comparison}
The six frozen classifications are:

\begin{itemize}
\item \textbf{Verified preserved:} The destination expresses the contracted property exactly or through a verified accepted equivalent with independent evidence.
\item \textbf{Observed partial:} Some required components survive, while another representable component is directly demonstrated to differ or be absent.
\item \textbf{Altered:} The feature remains present but an author-relevant value changes, such as alternative-text wording, language specificity, or title metadata.
\item \textbf{Confirmed lost:} The source property is valid and representable in PDF, but independent destination checks corroborate absence of the required representation.
\item \textbf{Unresolved equivalence:} Relevant structure or visible content exists, but the available evidence cannot establish cross-format equivalence; absence is not treated as loss.
\item \textbf{Measurement error:} The output is valid and relevant structure exists, but the available evidence cannot reliably resolve the contract.
\end{itemize}

Abstention is a scientific outcome: it prevents a parser limitation or representational ambiguity from becoming an unjustified pass or fail. The primary analysis uses the frozen comparator and taxonomy without post-result rule changes.

\section{V4 Experimental Design}
\label{sec:design}

\subsection{Feature families and variants}
The V4 benchmark selects 13 accessibility feature families that are both conversion-relevant and machine-observable across DOCX and tagged PDF. Table~\ref{tab:families} states the contract focus. Five controlled variants were created for every family to expose feature depth rather than inflate an arbitrary feature count. The variant patterns are retained in the source manifest and the final result matrix.

\begin{table*}[t]
\centering\scriptsize
\caption{Frozen V4 feature families and contract focus.}
\label{tab:families}
\begin{tabular}{p{0.16\textwidth}p{0.31\textwidth}p{0.45\textwidth}}
\toprule
Family & Source property & Contract focus \\
\midrule
Headings & Hierarchical paragraph semantics & Levels and sequence, not visual size \\
Alternative text & Author-written image description & Exact description associated with the image \\
Lists & Ordered/unordered list structure & Type, nesting, and list-item relationships \\
Basic tables & Table and header structure & Table, row, cell, and header roles \\
Document language & Primary document language & Document-level language value \\
Inline language & Run-level language override & Non-default language span inside an otherwise default-language document \\
Decorative images & Decorative-state semantics & Meaningful versus explicitly decorative image state \\
Hyperlinks & Link semantics & Link existence, visible name, and URI relation \\
Document metadata & Title and metadata identity & Visible title and document metadata values \\
Complex tables & Multi-level table semantics & Spans, scopes, and header associations \\
Footnotes & Note association & Reference-to-note relation and note body \\
Equations & Native mathematical structure & Formula representation beyond visible glyphs \\
Captions & Figure/caption relation & Caption text and association with the figure \\
\bottomrule
\end{tabular}
\end{table*}

The reading-order family was considered but excluded from the primary 13-family denominator because a defensible automated source-to-PDF oracle was not established. This is a measurement-boundary decision, not a claim that reading order is unimportant.

\subsection{Corpus and experimental units}
The 65 atomic DOCX fixtures are controlled and purposive. Each fixture isolates the target property while retaining enough surrounding document structure to exercise a real conversion. The primary unit is one fixture under one converter, yielding 65 $\times$ 2 = 130 cases. The corpus is not a random sample and does not support prevalence or inferential statistics.

\subsection{Supplementary Microsoft Word arm}
LibreOffice and Google Docs were the two pipelines in the prospectively frozen primary experiment. After that experiment was complete, Microsoft Word desktop's native GUI PDF export was recovered and its route was frozen before Word accessibility analysis. The supplementary condition used the same 65 frozen DOCX fixtures, source contracts, destination extraction, frozen comparator, and classification taxonomy. Word exported genuine tagged PDFs through \textit{File} $\rightarrow$ \textit{Export} $\rightarrow$ \textit{Create PDF/XPS}; no source editing, print-to-PDF route, or post-processing was used. Word was not retroactively added to the prospectively frozen two-pipeline denominator: it contributes 65 supplementary observations, while the complete archive contains 195 atomic fixture--pipeline outcomes.

The Word GUI route was used because COM activation remained unavailable under the recorded Windows session (0x80070520). This is an automation-interface limitation, not a substitution of a different converter. The Word protocol records the route, settings, input hashes, output provenance, and retry policy separately from the primary V4 protocol.

\subsection{Prospective freeze and holdout}
The V4 protocol, source contracts, fixture manifest, comparator taxonomy, decision rules, and holdout membership were frozen before final converter outputs were analyzed. The V4 freeze is preserved at commit \texttt{41a7d3f8b84e98acbdafe14070f2ebe7ea6a59e1}; the final hardening baseline records the pre-result hashes. V3 remains historical provenance and was not overwritten or used as a substitute for V4 outputs.

The 26-case holdout was pre-specified and kept separate from calibration and primary decision-rule development. It is a within-family/within-variant holdout: variants D and E are held out for every family, yielding one held-out case for each of the 13 families per primary pipeline. It is not a fully independent natural corpus, because all holdout documents are project-created members of the same controlled family and variant design. Holdout outcomes are reported separately rather than pooled with the 78 non-holdout atomic cases.

\subsection{Calibration and mutation testing}
The 78 calibration controls cover positive, negative, alternate-representation, and abstention situations. They yielded 60 correct decisions, zero false positives, five conservative false negatives, 13 correct abstentions, and zero unsupported confident classifications. The false negatives involve evidence boundaries for decorative-image alternatives, hyperlink name relations, and complex-table associations; expected labels were not changed after the converter results.

The source-oracle mutation audit applied 39 source mutations and detected all 39. This tests sensitivity to deliberate source perturbations; it is not an external validity or general mutation score for the whole benchmark.

\subsection{Realistic and repeatability arms}
The primary realistic arm contains 10 project-created realistic/integrated multi-feature documents, 78 source property contracts per pipeline, and 20 LibreOffice/Google PDFs. The supplementary Word route produced 10 additional realistic PDFs and the same 78 contracts were evaluated with the frozen procedure. These documents are not external real-world evidence and are not merged with the 130-case primary atomic denominator.

The repeatability arm is separate from the holdout. Six difficult-feature variants were converted three times per pipeline from identical source documents. LibreOffice and Google Docs each therefore have 18 repeated conversions; Word adds 18 supplementary conversions. The repeated classification and normalized semantic evidence are the stable units: Word was stable 6/6 on both, although its PDFs were byte-identical for 0/6 selected variants.

\section{Results}
\label{sec:results}

\subsection{Overall accounting}
Table~\ref{tab:overall} reports the complete primary atomic accounting, with the post-freeze Word arm shown separately as supplementary. No global score or weighted ranking is computed.

\begin{table}[t]
\centering\footnotesize
\caption{Primary atomic V4 outcomes for LibreOffice and Google Docs (P), with Microsoft Word reported separately as a supplementary condition (S). Each column contains 65 cases.}
\label{tab:overall}
\begin{tabular}{lrrr}
\toprule
Classification & LibreOffice (P) & Google Docs (P) & Word (S) \\
\midrule
\verified & 38 & 11 & 16 \\
\partialresult & 11 & 32 & 26 \\
\alteredresult & 9 & 11 & 16 \\
\lostresult & 1 & 6 & 2 \\
\unresolvedresult & 1 & 5 & 0 \\
\measurementresult & 5 & 0 & 5 \\
\bottomrule
\end{tabular}
\end{table}

The verified-preserved counts differ across these recorded conditions, while the three pipelines exhibit different mixtures of partial, altered, lost, unresolved, and measurement-error outcomes. These aggregate counts are not a product ranking: the corpus is controlled and purposive, the categories are not a weighted utility score, and Word's 65 observations are supplementary rather than directly comparable prospective treatment data.

\subsection{Feature-level matrix}
Table~\ref{tab:matrix} exposes the five variants within every family. Entries are in frozen source-manifest order and use V = verified preserved, P = observed partial, A = altered, L = confirmed lost, U = unresolved equivalence, and M = measurement error. The matrix is the central result because it preserves heterogeneity hidden by aggregate counts.

\begin{table*}[t]
\centering\scriptsize
\caption{V4 feature-family matrix across five variants. LibreOffice and Google Docs are primary (P); Word is supplementary (S). Entries are in frozen variant order.}
\label{tab:matrix}
\begin{tabular}{p{0.22\textwidth}p{0.25\textwidth}p{0.25\textwidth}p{0.25\textwidth}}
\toprule
Feature family & LibreOffice (P) & Google Docs (P) & Word (S) \\
\midrule
Headings & V/V/V/V/V & P/P/P/P/P & P/P/P/P/P \\
Alternative text & A/A/A/A/A & V/V/V/V/V & A/A/A/A/A \\
Lists & V/V/V/P/P & P/P/P/P/P & V/P/V/P/P \\
Basic tables & V/V/V/V/V & V/V/V/V/V & P/P/P/P/P \\
Document language & V/V/V/V/V & A/A/A/A/A & A/A/A/A/A \\
Inline language & V/V/V/V/A & L/L/L/L/L & A/A/A/A/A \\
Decorative images & P/A/P/P/P & P/V/P/P/P & P/A/P/P/P \\
Hyperlinks & P/P/A/A/L & P/P/P/P/A & P/P/P/P/L \\
Document metadata & V/V/V/V/V & A/A/A/A/A & V/V/V/V/V \\
Complex tables & P/P/V/V/P & P/P/P/P/P & P/P/P/P/P \\
Footnotes & M/M/M/M/M & U/U/U/U/U & M/M/M/M/M \\
Equations & V/V/V/V/V & P/P/P/P/P & V/V/V/V/V \\
Captions & V/V/U/V/V & P/P/L/P/P & V/V/L/V/V \\
\bottomrule
\end{tabular}
\end{table*}

The matrix shows several distinct behaviors. LibreOffice preserved headings, basic tables, document language, equations, and most inline-language variants, but altered all alternative-text variants. Google Docs preserved all alternative-text variants and basic tables, but all inline-language variants were confirmed lost. Word preserved all document-metadata and equation variants, while its headings and basic-table variants were partial and its alternative-text variants were altered. Footnotes remained measurement-bound in LibreOffice and Word and unresolved in Google Docs. These labels concern the contracted machine-observable semantics, not visible-text survival alone.

\subsection{Holdout, calibration, and reliability checks}
Table~\ref{tab:checks} separates the prospective holdout and diagnostic arms from the primary matrix.

\begin{table*}[t]
\centering\footnotesize
\caption{V4 validation and secondary-arm summary.}
\label{tab:checks}
\begin{tabular}{p{0.22\textwidth}p{0.24\textwidth}p{0.24\textwidth}p{0.24\textwidth}}
\toprule
Arm & LibreOffice (P) & Google Docs (P) & Word (S) \\
\midrule
26-case holdout & 14 preserved, 5 partial, 4 altered, 1 lost, 2 measurement errors & 4 preserved, 13 partial, 5 altered, 2 lost, 2 unresolved & 6 preserved, 11 partial, 6 altered, 1 lost, 2 measurement errors \\
Calibration & 78 controls: 60 correct, 0 false positives, 5 conservative false negatives, 13 correct abstentions & Same diagnostic set & Same diagnostic set \\
Source-oracle mutation & 39/39 deliberate source mutations detected & Same source audit & Same source audit \\
Repeatability & 6 variants $\times$ 3 runs; classifications stable 6/6; normalized semantics stable 6/6; byte identity 0/6 & 6 variants $\times$ 3 runs; classifications stable 6/6; normalized semantics stable 6/6; byte identity 6/6 & 6 variants $\times$ 3 runs; classifications stable 6/6; normalized semantics stable 6/6; byte identity 0/6 \\
Realistic arm & 10 project-created documents, 78 contracts, 20 PDFs total & Same primary arm & 10 supplementary PDFs; 78 contracts evaluated \\
\bottomrule
\end{tabular}
\end{table*}

The supplementary Word arm produced 16 verified-preserved, 26 observed-partial, 16 altered, 2 confirmed-lost, no unresolved-equivalence, and 5 measurement-error outcomes across the same 65 atomic fixtures. Because Word was evaluated after the prospectively frozen primary experiment, these 65 outcomes are reported separately rather than pooled into the 130-case primary denominator and are not directly comparable as prospective treatment data. The calibration false negatives are conservative evidence-boundary cases, not false-positive comparator behavior. In the repeatability audit, normalized semantic evidence and classifications remained stable 6/6 for all three pipelines. Google PDFs were byte-identical for all six selected variants; Word PDFs were byte-identical for none of the six. LibreOffice byte differences were attributable to metadata such as creation time; the normalized semantic record remained stable. Raw extraction identity is not claimed as a scientific stability unit because records contain route-specific paths and process-local object addresses.

\subsection{Confirmed losses, alterations, and unresolved cases}
The Google inline-language rows are confirmed lost under the contract: the source contains a non-default run-level language value, while independent destination checks find no required inline language span or equivalent marker. Visible text remains, but visible text is not the source language semantic. The Google footnote rows likewise retain visible reference and note text but lack independently recoverable machine-verifiable association structure.

The nine LibreOffice altered cases include the five alternative-text variants, where the destination retains related image description content but changes the author-provided value; the Google altered cases include document language, metadata, and one hyperlink result. The five Google unresolved footnote rows and the one LibreOffice unresolved caption variant remain unresolved rather than being forced into loss or preservation. LibreOffice's five footnote measurement errors retain relevant note structure, but the evidence cannot establish the body/reference relation reliably.

\subsection{Realistic documents}
The realistic arm produced 53 preserved, 5 partial, 15 altered, and 5 measurement-error observations for LibreOffice, and 20 preserved, 33 partial, 10 altered, and 15 confirmed-lost observations for Google Docs. These 156 property--pipeline observations show that the contracts remain usable when features coexist, but they are project-created integrated evidence, not a natural external corpus and not a prevalence estimate.

The supplementary Word realistic arm adds 10 genuine Word PDFs and 78 contract observations: 23 verified-preserved, 25 observed-partial, 25 altered, and 5 measurement-error outcomes. It is reported separately from the primary realistic arm.

\section{Discussion}
\label{sec:discussion}

\subsection{Destination accessibility and preservation are different questions}
The central methodological result is that destination accessibility and source-semantic preservation should be reported separately. A tagged PDF can contain headings, figures, tables, and visible text while changing a source language span, title value, alternative-text string, or association between a footnote and its reference. Conversely, a destination structure can be present while the available evidence is insufficient to establish that it represents the source contract. The unresolved and measurement-error categories prevent both situations from being hidden inside a single pass/fail label.

\subsection{Testing implications}
For a software-testing audience, the protocol makes the testing chain explicit: contracts are feature-specific oracles, DOCX files are controlled test inputs, each converter is a system under test, PDF extraction produces observations, and precedence rules make the oracle decision. Calibration controls expose alternate encodings that cannot safely be accepted by the current comparator. Such cases become actionable unresolved results rather than silent false failures or false passes.

The mutation audit provides evidence that the source oracle detects deliberate changes, while the holdout provides a prospective check on the frozen decision procedure. These are complementary: mutation detection tests sensitivity to source perturbation; holdout separation tests that the primary corpus was not used as an informal tuning set. Neither establishes external validity or complete fault-model coverage.

\subsection{Pipeline-specific observations}
The feature matrix shows no monotonic converter pattern. Google Docs preserved the author-written alternative-text strings in all five variants while LibreOffice and Word altered them; LibreOffice preserved more explicit structural evidence for several list, language, equation, and caption cases; Google Docs confirmed loss of every tested inline-language variant, while Word altered those values. The supplementary Word arm also preserved all tested document-metadata and equation variants, but showed partial headings and basic-table behavior. The Google condition includes DOCX import, an internal document representation, and PDF generation, while Word used native desktop export; the evidence establishes end-to-end behavior under each recorded workflow, not the location of a defect within a converter.

Repeatability strengthens the operational claim for the recorded conditions. The six selected classifications and normalized semantic evidence were stable across three repeated conversions for each of the three pipelines (6/6 per pipeline). Word's PDFs were not byte-identical across the six selected variants, so byte identity and semantic-result reproducibility remain distinct claims. This does not make a cloud service timeless, nor does it turn byte identity into a requirement for semantic repeatability.

\subsection{Engineering and accessibility implications}
The protocol can support regression testing in authoring, publishing, and document-conversion pipelines. A converter adapter can be evaluated against the same frozen sources and contracts, while unresolved cases identify where parser or representational work is needed. The practical implication is not that every export must be rejected or that one product should always be selected. It is that destination-only checks should be complemented by checks of selected author-supplied properties when preservation matters.

The classifications intentionally do not encode user impact. An altered description may remain understandable, an observed partial structure may retain substantial accessibility, and a confirmed machine-level loss may have different consequences for different tasks. Connecting structural preservation to task performance requires a separate user-centered study.

\subsection{Contribution boundary}
A11yPreserve is a source-contract conformance protocol, benchmark, comparison implementation, and empirical study. It is not a complete PDF accessibility validator, a user study, a population survey, a universal converter ranking, or a claim that PDF export alone caused every observed difference. The prospectively frozen primary evidence covers two named pipeline conditions; the recovered Word route supplies a separately frozen supplementary third-pipeline observation set.

\section{V5 Hardening and Profile Audit}
\label{sec:v5-hardening}

V5 adds independent checks around the frozen V4 experiment without silently changing its primary classifications. The hardening layer is additive: it records chronology, source and destination fault sensitivity, holdout membership, profile settings, and implementation disagreement so that a reviewer can distinguish stronger evidence from a changed result denominator.

\subsection{Profile interpretation and condition comparison}
\label{sec:v5-profile}

The V4 LibreOffice condition used \texttt{UseTaggedPDF=true}, \texttt{PDFUACompliance=true}, and \texttt{SelectPdfVersion=1}. The V5 profile audit documents that the latter selects PDF/A-1b rather than an ordinary default PDF export. V5 condition B repeated all 65 LibreOffice fixtures with \texttt{SelectPdfVersion=0}, while retaining the tagging and PDF/UA settings. All 65 condition-B conversions produced valid outputs; all were tagged, and the raw PDF/A marker count was zero. The machine-derived classifier produced zero semantic classification changes relative to V4. This is evidence of no material effect on the tested semantic classifier, not evidence of byte identity, complete PDF/UA conformance, or equivalence of every PDF serialization detail.

\subsection{Holdout accounting and fault sensitivity}
\label{sec:v5-holdout}

The V5 holdout audit derives fixture membership from the frozen manifests rather than from prose counts. The 130 primary fixture--pipeline cases contain 52 fixture-level holdout cases (26 per primary pipeline) and 78 non-holdout cases (39 per pipeline). The holdout is not pooled with calibration controls, and the supplementary Word arm remains outside the primary denominator.

The source mutation audit detected 39 of 39 deliberate source mutations. Destination-side mutation testing applied the family-specific operators across all 65 atomic LibreOffice PDFs, spanning five variants per family and multiple structure/value locations. Fifty-nine attempts produced valid destination mutants; 46 changed the frozen classification, 13 survived without a classification change, six attempts were unsupported, and two valid mutants produced measurement errors. Sixty-five no-op controls retained their baseline classifications with zero false-positive controls. These tests establish sensitivity and expose blind spots for the recorded fault models, not exhaustive coverage of all converter or PDF-parser faults.

\subsection{Independent diagnostic cross-check and calibration}
\label{sec:v5-validation}

An independent lightweight raw/PyMuPDF implementation inspected 26 rows, one per feature family per primary pipeline. It agreed with the original classification on 8 rows and disagreed on 18. The disagreement is informative: the second implementation intentionally lacked some exact-value and association logic, so it frequently abstained or could not establish the richer contract. It is therefore reported as an independent diagnostic cross-check, not as an external adjudicated ground truth and not as a replacement for the frozen comparator.

The 78-case calibration audit produced 60 correct decisions, zero false positives, five conservative false negatives, and 13 correct abstentions. The calibration result supports the claimed conservative behavior for the tested controls; it is not a sensitivity, specificity, or population-performance estimate.

\subsection{What V5 establishes}
\label{sec:v5-interpretation}

Taken together, the V5 checks strengthen four claims. First, the primary accounting is prospectively separable into holdout and non-holdout cases. Second, the comparator detects the tested source and destination perturbations. Third, the main result is not an artifact of the documented LibreOffice PDF/A profile interpretation. Fourth, a second technical implementation exposes where exact cross-format equivalence is difficult rather than hiding that uncertainty.

The hardening does not convert the study into a global accessibility benchmark. The atomic corpus remains controlled and purposive; the realistic documents are project-created; Google Docs remains a cloud condition whose backend cannot be pinned; the Word arm is supplementary; reading order remains outside the tested attribute set; and no disabled-user or screen-reader task study was performed. These are properties of the research question and evidence boundary, not packaging defects. The final claim is consequently precise: A11yPreserve provides an auditable source-contract testing methodology for selected machine-verifiable accessibility semantics across the recorded DOCX-to-PDF conditions.

\section{Threats to Validity}
\label{sec:threats}

\subsection{Construct validity}
Cross-format equivalence is not identical to string equality or tag presence. PDF permits multiple encodings for related semantics, and the current extractor cannot establish every possible equivalent representation. The comparator therefore uses feature-specific contracts and abstains when evidence is insufficient. The five calibration false negatives, six unresolved atomic outcomes, five LibreOffice footnote measurement errors, and deferred reading-order family make this boundary explicit. The study does not claim complete PDF accessibility or standards conformance.

\subsection{Internal validity}
The source oracle and comparator were implemented and reviewed by the same research team, so their agreement is not external validation. The 39/39 mutation result tests deliberate source changes but does not prove oracle completeness. The conversion records preserve source/output hashes, routes, timestamps, and producer observations, and the V4 methodology was frozen before final result inspection. The holdout remained separate from calibration and primary analysis. These controls reduce but do not eliminate implementation and process error.

\subsection{External validity}
The 65 atomic fixtures are controlled, purposive, and project-generated. Each primary manifest has 75 records, comprising those 65 atomic fixtures plus 10 project-created realistic integrated documents. The 26-case-per-pipeline holdout is within-family/within-variant, not a fully independent natural corpus. The realistic arm is also project-created and integrated; it is not a natural real-world corpus. The prospectively frozen primary comparison contains two pipelines. Microsoft Word was recovered only after primary-result completion and was therefore evaluated as a supplementary third pipeline rather than retroactively incorporated into the primary design. Results do not generalize to all DOCX documents, converters, software versions, or accessibility properties.

\subsection{Reliability and environment dependence}
LibreOffice was tested under a recorded local version and export configuration. Google Docs is a cloud service whose backend version cannot be independently pinned. Word used native desktop GUI export because COM automation was unavailable in the local Windows session; the route, settings, source hashes, and output provenance were frozen and recorded, but reviewers may not reproduce the COM failure or GUI condition identically in every Windows environment. The repeatability arm shows stable classifications for six selected variants under the recorded environment, but a later service or application update can change behavior. Output byte identity is not required for semantic stability; the final audit treats normalized evidence and classification identity separately from process-local extraction metadata.

\subsection{User validity and ethics}
No disabled-user study, screen-reader task study, perceived-usability study, or severity assessment was performed. The results concern machine-verifiable preservation of selected source information, not task completion or user harm. The project used controlled and project-created documents and did not upload an external natural-document corpus. The anonymous replication package contains no credentials, browser state, or private account material.

\subsection{Interpretation}
Counts are descriptive controlled-case summaries. No prevalence estimate, inferential statistic, weighted accessibility score, or ecosystem-wide product ranking is justified. An altered or partial property may remain usable, an unresolved case may later prove equivalent under richer evidence, and a confirmed machine-level loss does not by itself establish human harm. These limitations define the valid scope of the conclusions rather than invalidating the source-grounded measurement question.

\section{Conclusion}
\label{sec:conclusion}

A11yPreserve provides a source-grounded way to test whether selected accessibility semantics survive DOCX-to-PDF conversion. The frozen V4 design combines feature contracts, raw OOXML verification, a separately implemented source oracle, real conversion pipelines, multiple destination evidence paths, conservative classifications, prospective holdout evaluation, calibration, mutation testing, repeatability, and a separate realistic-document arm.

The prospectively frozen primary evaluation contains 130 LibreOffice and Google Docs fixture--pipeline cases. Each primary manifest contains 75 records: 65 atomic fixtures and 10 project-created realistic integrated documents. LibreOffice produced 38 verified-preserved, 11 observed-partial, 9 altered, 1 confirmed-lost, 1 unresolved, and 5 measurement-error outcomes. Google Docs produced 11, 32, 11, 6, 5, and 0 respectively. A later full-corpus supplementary Microsoft Word evaluation adds 65 observations: 16 verified-preserved, 26 observed-partial, 16 altered, 2 confirmed-lost, no unresolved-equivalence, and 5 measurement-error outcomes. Six repeated variants were classification-stable and semantically stable 6/6 for each pipeline, although Word was byte-identical for 0/6. The 10 project-created realistic documents and 78 contracts per pipeline remain a separate applicability arm; Word adds 10 supplementary PDFs and 78 separately reported contract observations.

The evidence establishes condition-specific preservation behavior, not a global converter ranking or complete accessibility assessment. It also does not establish PDF/A or PDF/UA conformance, byte-identical serialization, or human usability. The strongest methodological lesson is that destination accessibility and preservation of known source semantics are different testing questions. A complete submission should therefore report the source contracts, destination evidence, uncertainty boundary, chronology of the two-pipeline primary design and the later Word supplement, and provenance together rather than infer preservation from tagged output alone.


\section*{AI Disclosure}
OpenAI ChatGPT and OpenAI Codex assisted with drafting and restructuring portions of the manuscript; with adversarial review; and with verification-script, submission-artifact, and package documentation. The assistance was at the prose, organization, critique, code-support, and packaging level. These systems were not used as experimental ground truth, independent human validation, or autonomous decision-makers for frozen classifications or final scientific conclusions. Human authors reviewed and verified AI-assisted material against the source artifacts, frozen evidence, analysis records, and cited literature.

\bibliographystyle{IEEEtran}
\bibliography{references}

\end{document}
